\section{Conclusion}

Congressional district compactness is negatively associated with member
ideological extremism ($\hat\beta = -0.071$ per 0.10 PP, $p < 0.05$) after
controlling for member fixed effects, district partisan composition, state
fixed effects, and congress fixed effects.
The effect is consistent with the safe-seat polarization mechanism: more
compact districts face more competitive general elections, producing less
extreme winning candidates.

The counterfactual projection suggests that replacing the five study states'
enacted gerrymanders with \textsc{bisect} algorithmic plans would reduce
mean absolute DW-NOMINATE by approximately 0.072 points --- a 13\% reduction.
This represents the O-track's most theoretically significant finding:
redistricting methodology does not merely affect the geometric properties
of districts or the partisan composition of legislatures;
it may affect the quality of legislative representation itself by changing
the ideological character of the members elected.

Two caveats deserve emphasis.
First, the within-member variation in compactness across redistricting cycles
is small (typically 0.05--0.15 PP), limiting statistical precision.
Second, the projection assumes the within-member slope applies to
between-plan comparisons; this extrapolation requires causal identification
that the present analysis cannot provide.
These limitations motivate a Phase~2 analysis using redistricting cycle
boundaries as instruments.
